% Copyright 2017 Rebecca Skinner
%
% This work is licensed under the Creative Commons
% Attribution-ShareAlike 4.0 International License. To view a copy of
% this license, visit http://creativecommons.org/licenses/by-sa/4.0/
% or send a letter to Creative Commons, PO Box 1866, Mountain View, CA
% 94042, USA.
\documentclass{beamer}

\title{A Brief Survey of Machine Learning in Haskell}
\author{Rebecca Skinner\\ \small{@cercerilla}}
\date{\today}

\mode<presentation> {\usetheme{metropolis}}

\usepackage[english]{babel}
\usepackage{times}
\usepackage[T1]{fontenc}
\usepackage{hyperref}
\usepackage{listings}
\usepackage{color}
\usepackage{amsmath}
\usepackage{csquotes}
\usepackage{verbatim}
\usepackage{fontspec}
\usepackage{pbox}

\definecolor{comment}{rgb}{145,175,188}
\definecolor{keyword}{rgb}{157,163,199}
\definecolor{string}{rgb}{155,204,174}

\lstset{ % add your own preferences
  basicstyle=\tiny,
  showspaces=false,
  showtabs=false,
  numbers=none,
  numbersep=5pt,
  showstringspaces=false,
  stringstyle=\color[rgb]{0.16, .47, 0},
  tabsize=1
}

\newcommand{\chref}[3] {
  {\color{#1} \href{#2}{\underline{#3}}}
}

\AtBeginSection[]{
  \begin{frame}
    \vfill
    \centering
    \begin{beamercolorbox}[sep=8pt,center,shadow=true,rounded=true]{title}
      \usebeamerfont{title}\insertsectionnumber \\ \insertsectionhead\par%
    \end{beamercolorbox}
    \vfill
  \end{frame}
}

\AtBeginSubsection[]{
  \begin{frame}
    \vfill
    \centering
    \begin{beamercolorbox}[sep=8pt,center,shadow=true,rounded=true]{title}
      \usebeamerfont{title}\insertsectionnumber.\insertsubsectionnumber\\\insertsubsectionhead\par%
    \end{beamercolorbox}
    \vfill
  \end{frame}
}

\begin{document}
\begin{frame}
  \titlepage{}
  \begin{center}
    \small{\chref{blue}{http://creativecommons.org/licenses/by-sa/4.0/}{LICENSE}}
  \end{center}
\end{frame}

\section{Introduction}

\begin{frame}
  \frametitle{A Bit About Haskell}
  Haskell is a strongly typed pure functional programming language.
  It's used in industry, as a research language, and for teaching.  It
  has broad name recognition, but is (somewhat unfairly) derided for
  being ``overly academic''.
\end{frame}

\begin{frame}
  \frametitle{What We'll Talk About}
  This talk is going to be a survey of haskell as a language for
  building applications that rely on data science and leverage machien
  learning.  We'll talk about why haskell is well suited to these
  types of applications, what tools exist, and where you might run
  into problems.
\end{frame}

\begin{frame}
  \frametitle{What We Won't Talk About}
  This talk is going to focus on introducing haskell to users who are
  already familiar with, or working on learning how to use machine
  learning for building their applications.  We won't be diving into
  the details of specific approaches.
\end{frame}

\section{The Current State of ML in Haskell}

\begin{frame}
  \frametitle{Haskell is Used for Machine Learning}
  There are a few companies that are using haskell as part of their
  data analysis and ML workloads.  It's not a major player in the
  space, but there is some support from a few big names:
  \begin{itemize}
  \item Galois: Machine Learning for Cyber Security
  \item Facebook: Building DSLs for Anti-Abuse Engines
  \item Target: Logistics and Consumer Behavior
  \item Takt (Starbucks): Rewards, Consumer Behavior
  \item Rackspace: Business Intelligence, Analytics, Support Automation
  \item Microsoft: R\&D
  \item HFT and Quants: Trading Algorithms
  \end{itemize}
\end{frame}

\begin{frame}
  \frametitle{DataHaskell}
  DataHaskell is a group dedicated to improving machine learning and
  data science stories in haskell:
  \url{http://www.datahaskell.org/}
\end{frame}

\begin{frame}
  \frametitle{Statistical Computing}
  Haskell has strong support for statistical computing through the
  haskell stats package as well as through support for the GNU
  Scientific Library (GSL).
\end{frame}

\begin{frame}
  \frametitle{Linear Algebra}
  Haskell has both native libraries for linear algebra, as well as
  lightweight wrappers around libraries like GSL, including libraries
  that offer support for GPGPU computations.
\end{frame}

\begin{frame}
  \frametitle{Hardware Accelerated ML}
  The most popular haskell compiler, GHC, supports x86 and x86-64
  systems, with growing support for ARM and PowerPC.\@  There are,
  however, numerous projects to allow haskell to build to, or
  integrate with, other architectures including nVidia GPUs and
  compiling haskell code directly to hardware description languages.
\end{frame}

\begin{frame}
  \frametitle{Interoperability with C}
  GHC has a mature and well supported FFI that allows it to interact
  with C libraries.  This means that haskell can easily support any
  machine learning and general purpose mathematical libraries that are
  written in C.
\end{frame}

\begin{frame}
  \frametitle{Library Support}
  There are three large and well documented haskell libraries that
  support high level out-of-the-box machine learning:
  \begin{itemize}
  \item HLearn
  \item Grenade
  \item haskell-tensorflow
  \end{itemize}
\end{frame}

\subsection{Haskell Machine Learning Libraries}

\begin{frame}
  \frametitle{Grenade}
  A dependently typed DSL in pure haskell written to support building
  and composing neural networks.  Actively developed.
\end{frame}


\begin{frame}
  \frametitle{Tensorflow Haskell}
  Provides a rich set of idiomatic libraries on top of libtensorflow.
  Actively developed, but only supports TensorFlow <= 1.3.
\end{frame}

\begin{frame}
  \frametitle{HLearn}
  Built to support research into homomorphic machine learning
  algorithms.  Pure haskell, with high performance.  Deprecated.
\end{frame}

\section{Where Haskell May Fall Short}

\begin{frame}
  \frametitle{Out Of Date Libraries}
  Of the three major libraries available to provide out-of-the-box ML
  capabilities in haskell, one is deprecated, and the only only
  supports an outdated version of TensorFlow.  Most new work being
  done in the field is not well documented.
\end{frame}

\begin{frame}
  \frametitle{Missing Libraries}
  The lack of public open source work in data analysis and machine
  learning may mean that there are fewer libraries available for these
  tasks compared to other languages.
\end{frame}

\begin{frame}
  \frametitle{Proprietary Work}
  Many companies are doing ML in haskell, and actively hiring, but
  much of the code being developed is proprietary.  This means that it
  can be difficult to get started without a team dedicated to building
  tooling from the ground up.
\end{frame}

\begin{frame}
  \frametitle{Performance Concerns}
  Although haskell is capable of high performance, it can be difficult
  to achieve in practice.  Lazy evaluation can lead to unexpected
  runtimes and much higher than expected memory utilization.
  Integrating with code running on GPUs, FPGAs, and ASICs can be
  difficult if you're not already familiar with the GHC internals.
\end{frame}


\begin{frame}[fragile]
  \begin{center}
    \includegraphics[width=.85\paperwidth]{images/typesig.png}
  \end{center}
\end{frame}

\begin{frame}
  \frametitle{Cognative Overhead}
  Many of the libraries available impose a great deal of rigor into
  how they represent the ML models available.  This can lead to a lot
  of additional cognitive overhead when exploring a problem space if
  you're unaccustomed to working under those constraints.
\end{frame}

\section{Reasons to Consider Haskell}

\begin{frame}
  \frametitle{Why Use Haskell for ML?}
  In spite of the overall immaturity of ML in the haskell ecosystem, there are several compelling reasons to look at using haskell for data science and machine learning applications.  These come from three major areas:
  \begin{itemize}
    \item Expressiveness
    \item Performance
    \item Correctness
  \end{itemize}
\end{frame}

\subsection{Expressiveness}

\begin{frame}[fragile]
  \begin{center}
    \includegraphics[height=.85\paperheight]{images/scream.png}
  \end{center}
\end{frame}

\begin{frame}
  \frametitle{What Is Expressiveness?}
  Expressiveness speaks to the ability of a user to clearly and
  concisely represent their thoughts in a language, with a minimum
  amount of extraneous boilerplate.  Because machine learning and data
  science are intricately tied to underlying mathematical notions of
  computation, the syntax and semantics of haskell are particularly
  well suited to expressing problems in these domains.
\end{frame}

\begin{frame}
  \frametitle{Composability}
  Composability is really about how easy it is to build complex things
  from smaller pieces.  By convention, haskell libraries, including
  the machine learning libraries we've discussed, focus on
  composability.
\end{frame}

\begin{frame}
  \frametitle{DSLs}
  Domain Specific Languages (DSLs) are a way of creating a small
  language, or language-like API for a library, that help you express
  your problem.  The Grenade and HLearn libraries both focus heavily
  on implementing DSLs for machine learning problems, and haskell
  lends itself very well to this approach.

  By providing an easy way to implement DSLs, you can wrap ML and data
  analysis capabilities of your application in an easy-to-use frontend
  to guard against misuse.
\end{frame}

\subsection{Performance}

\begin{frame}[fragile]
  \begin{center}
    \includegraphics[width=.85\paperwidth]{images/santa.jpg}
  \end{center}
\end{frame}

\begin{frame}
  \frametitle{Performance vs Python and Scala}
  Haskell is a compiled language with a long history of work on
  performance optimizations that help make code performant.  In
  typical scenarios, haskell runs about as fast as C++ or Java, with a
  somewhat higher memory footprint, and much faster than pure python
  code with a somewhat smaller memory footprint.
\end{frame}

\begin{frame}
  \frametitle{Performance vs NumPy}
  For libraries like NumPy that use highly optimized native code, pure
  haskell code will tend to be slightly slower.  Use of the FFI can
  mitigate this at the cost of some additional code and optimization
  complexity.
\end{frame}

\begin{frame}
  \frametitle{Parallelism}
  Haskell has a wide range of primatives and library suppor for
  parallel processing.  This can allow haskell applications to easily
  leverage available hardware.
\end{frame}

\begin{frame}
  \frametitle{Concurrency}
  Haskell has native support for several different concurrency models.
  Of particular utility is the Haxl library, which greatly simplifies
  the developemnt of applications that rely on asynchronous data
  sources.  Combined with lazy evaluation and composibility as a first
  class citizen, this can lead to efficiency gains when dealing with
  expensive computations.
\end{frame}

\begin{frame}
  \frametitle{Profiling}
  Tools like ThreadScope can help provide a visual way to debug
  performance and concurrency bugs in performance critical
  applications.
\end{frame}

\subsection{Correctness: The Benefits of Purity Type Safety}

\begin{frame}[fragile]
  \begin{center}
    \includegraphics[height=.85\paperheight]{images/ml.png}
  \end{center}
\end{frame}

\begin{frame}
  \frametitle{Referential Transparency}
  Haskell is a pure language, meaning that haskell functions are
  functions in mathematical sense.  Thanks to this, we can make
  assumptions about the behavior of code that we're reading that we
  might not be able to make about code written in other languages.

  Knowing that our functions are pure allows us to better create
  mental models of how our software is executing, removing uncertainty
  about the behaivor of critical parts of the application.
\end{frame}


\begin{frame}
  \frametitle{Algebraic Reasoning}
  Strongly related to the notion of purity is algebraic reasoning.
  Haskell idiomatically provides algebraic laws for entitites defined
  in the language.  When using built-in algebraic structures like
  monoids, monads, functors, semigroups, and rings we can reason about
  the structure of the code using the same tools that we would use to
  reason about them as mathematical structures.

  This is highly beneficial when developing novel or critical tools
  for analysis and machine learning, since the code that implements
  the transformations and models more closely aligns to the
  theoretical work we've done.
\end{frame}

\begin{frame}[fragile]
  \frametitle{Monad Laws}
    \begin{center}
    \includegraphics[width=.90\paperwidth]{images/monad-laws.png}
  \end{center}
\end{frame}

\begin{frame}
  \frametitle{Phantom Types}
  Phantom types allow you to track the provenance of the data as part
  of it's value.  By keeping track of the source of the data, you can
  make good decisions about how to treat it later in your application.
  For machine learning applications this can be particularly useful as
  you can know how reliable a given set of data may be based on where
  it originated.
\end{frame}

\begin{frame}[fragile]
  \frametitle{Phantom Types}
    \begin{center}
    \includegraphics[width=.85\paperwidth]{images/casper.jpg}
  \end{center}
\end{frame}

\begin{frame}
  \frametitle{Defining Data}
  Haskell's type system does more than allowing algebraic reasoning.
  By allowing us to give specific names to differnt kinds of data, we
  can ensure that we remember, and communicate to others, what the
  data should be.
\end{frame}

\begin{frame}
  \frametitle{Dependent Types}
  Dependent types allow the type of a variable to depend on it's
  value.  This means that we can write expressions that will, for
  example, a vector has a certain size, or that the dimensions of two
  matrices are correct when multiplying them.
\end{frame}

\begin{frame}[fragile]
  \frametitle{Dependent on Data}
    \begin{center}
    \includegraphics[width=.85\paperwidth]{images/dependent-data.jpg}
  \end{center}
\end{frame}

\section{Conclusion}

\begin{frame}
  \frametitle{Clear Advantages: Sometimes}
  Haskell offers some clear advantages in some circumstances, but the
  overhead of realizing them means that it might not be a good choice
  for exploratory projects, prototypes, or as a language for someone
  trying to learn more about ML
\end{frame}

\begin{frame}
  \frametitle{Great for Advanced Users}
  While haskell doesn't have a lot of batteries included, it's
  advantages as an expressive and performant language make it an ideal
  choice for users who are already pushing the boundries of what can
  be done with off-the-shelf libraries and need to implement their own
  custom solutions.
\end{frame}

\begin{frame}
  \frametitle{A Bad Way To Learn}
  While there are a few batteries included libraries available, they
  are written with an audience who is already deeply familiar with
  machine learning and data science.  These libraries make no attempt
  to be an easy introduction into the fundamental concepts.  Because
  of this, haskell would make a poor choice of language for someone
  wanting to start learning more about data science or machine
  learning.
\end{frame}

\section{Questions?}


\end{document}
